% ====================================================================
% filterResponse.tex
% --------------------------------------------------------------------
% Magnitude response of the filter at its nominal values, drawn with
% pgfplots from the design equation (f0 = 995.9 Hz, Q = 0.7071).
% ====================================================================

\begin{figure}[htbp]
\centering
\begin{tikzpicture}
\begin{semilogxaxis}[
  width=0.92\linewidth,
  height=7.5cm,
  grid=both,
  minor grid style={gray!20},
  major grid style={gray!45},
  xmin=10, xmax=20000,
  ymin=-60, ymax=5,
  ytick={-60,-50,-40,-30,-20,-10,0},
  xlabel={Frequency (Hz)},
  ylabel={Gain (dB)},
  legend pos=south west,
  legend cell align=left,
]
  \addplot[domain=10:20000, samples=300, very thick, linkblue]
    {-10*log10((1 - (x/995.9)^2)^2 + (x/(995.9*0.7071))^2)};
  \addlegendentry{Filter response}

  % Stopband requirement: 35 dB at 9.5 kHz
  \addplot[only marks, mark=*, mark size=2.5pt, red] coordinates {(9500,-35)};
  \addlegendentry{Requirement R5}

  % Top of the signal band
  \draw[dashed, gray] (axis cs:500,-60) -- (axis cs:500,5);
  \node[font=\footnotesize, anchor=south east, rotate=90]
    at (axis cs:500,-58) {top of signal band};

  \node[font=\footnotesize, anchor=west]
    at (axis cs:1100,-3) {\qty{-3}{\decibel} at \qty{996}{\hertz}};
  \node[font=\footnotesize, anchor=east]
    at (axis cs:9000,-41) {\qty{-39.2}{\decibel}};
\end{semilogxaxis}
\end{tikzpicture}
\caption{Calculated magnitude response at the nominal component values. The
curve passes \qty{4.2}{\decibel} below the stopband requirement.}
\label{fig:filterResponse}
\end{figure}
